\begin{questionS}
Write a concurrent program to find a maximum element in a large array~|a| of
|Int|s, using the bag-of-tasks pattern.  Each task should correspond to
finding the maximum in some segment of the array.  Give your code the
following signature:
%
\begin{scala}
class ArrayMax(a: Array[Int], numWorkers: Int, taskSize: Int){
  require(a.nonEmpty && numWorkers > 0 && taskSize > 0)

  def apply(): Int = ...
}
\end{scala}
%
where |numWorkers| is the number of workers, and |taskSize| is the number of
elements in each task; the |apply| function returns the maximum. 

Test your code.  

Carry out some informal experiments to find a sensible size for each task.  
\end{questionS}

%%%%%%%%%%

\begin{answerS}
My answer is below.
%
\begin{scala}
class ArrayMax(a: Array[Int], numWorkers: Int, taskSize: Int){
  require(a.nonEmpty && numWorkers > 0 && taskSize > 0)

  /** A £Task (l,r)£ represents the task of calculating the maximum of £a[l..r)£. */
  type Task = (Int, Int)

  /** Channel from the controller to workers. */
  private val toWorkers = new BuffChan[Task](numWorkers)

  /** Channel from the workers to the controller. */
  private val toController = new BuffChan[Int](numWorkers)

  /** A worker thread. */
  private def worker = thread("worker"){
    var myMax = Int.MinValue
    repeat{
      val (l,r) = toWorkers?()
      for(i <- l until r) myMax = myMax max a(i)
      }
    toController!myMax
  }

  /** The maximum value so far. */
  private var theMax = Int.MinValue

  /** The controller. */
  private def controller = thread{
    // Distribute tasks.
    var l = 0; val n = a.length
    while(l < n){ val r = (l+taskSize) min n; toWorkers!(l,r); l = r } 
    toWorkers.endOfStream()
    // Receive workers' maxima, and combine, storing the result in £theMax£
    for(_ <- 0 until numWorkers) theMax = theMax max toController?()
  }

  /** The system. */
  private def system = 
    controller || (|| (for (_ <- 0 until numWorkers) yield worker))

  /** Find the maximum. */
  def apply(): Int = { run(system); theMax }
}
\end{scala}

I tested my code using the standard technique of generating random data, and
comparing the result of the concurrent program with that from sequential
code.  The following function performs a single test; the main testing program
calls this many times.
% 
\begin{scala}
  def doTest() = {
    val a = Array.fill(1+Random.nextInt(1000))(Random.nextInt(Int.MaxValue))
    val numWorkers = 1+Random.nextInt(20); val taskSize = 1+Random.nextInt(100)
    val max = new ArrayMax(a, numWorkers, taskSize)()
    assert(max == a.max, a.mkString(",")+s"\nmax = $max")
  }
\end{scala}

I ran some informal experiments using an input array of size 10,000,000, and
eight workers, repeated 100 times.  I found that a task size of about 40,000
was best, taking around 9ms per array.  Using larger tasks (while still
considerably smaller than the whole array) gave a similar running time.
However, with a task size of 10,000, the running time was nearly twice as
large, as the channel communication became a bottleneck.
\end{answerS}
